\chapter{Arm Pain}

\section*{Definition}
Arm pain is pain or discomfort anywhere from the shoulder to the hand. It may arise from muscles, tendons, joints, bones, nerves, blood vessels, or referred pain from the neck, chest, or another area.

\section*{When to See a Doctor}

\begin{itemize}
 \item Arm pain with chest pressure or pain, shortness of breath, sweating, nausea, or fainting -- call emergency services because heart attack symptoms may spread to either arm
 \item A cold, pale, blue, or pulseless arm; sudden severe swelling; or sudden weakness or numbness
 \item Severe pain after trauma, deformity, inability to move the arm, fever with a hot swollen joint, or rapidly worsening symptoms
\end{itemize}

\section*{How It Develops}
Arm pain may be nociceptive from a strain, fracture, tendon, or joint; neuropathic from a nerve root or peripheral nerve; vascular from reduced blood flow or venous obstruction; or referred from the neck, chest, or shoulder. Arm pain from heart ischemia can occur on either side and may occur without obvious chest pain.

\section*{Causes}
\begin{itemize}
 \item Cervical radiculopathy (herniated disc, foraminal stenosis)
 \item Rotator cuff tendinopathy or tear
 \item Fracture (humerus, radius, ulna)
 \item Lateral epicondylitis (tennis elbow)
 \item Medial epicondylitis (golfer's elbow)
 \item Carpal tunnel syndrome
 \item Thoracic outlet syndrome
 \item Bursitis (olecranon, subacromial)
 \item Referred cardiac pain (angina or myocardial infarction, often but not always involving the left arm)
 \item Peripheral neuropathy (diabetic, alcoholic)
\end{itemize}

\section*{Related Symptoms}
\begin{itemize}
 \item Shoulder pain
 \item Neck pain
 \item Hand numbness
 \item Weakness
 \item Swelling
\end{itemize}


\section*{Medical Evaluation}
\begin{itemize}
 \item The clinician will examine the neck, shoulder, elbow, forearm, wrist, and hand, including strength, sensation, reflexes, pulses, range of motion, and selected provocative tests.
 \item X-rays are used when fracture, dislocation, arthritis, or another bone problem is suspected; the appropriate views depend on the injured area.
 \item MRI or other advanced imaging is reserved for selected suspected tendon, nerve, joint, vascular, or occult-bone problems. For persistent cervical radiculopathy, MRI of the cervical spine without contrast is often appropriate; CT with contrast is not a routine test.
 \item Nerve conduction studies and electromyography (EMG) are arranged when carpal tunnel syndrome, cubital tunnel syndrome, or cervical radiculopathy is considered.
\end{itemize}

\section*{Treatment Options}
\begin{itemize}
 \item Physiotherapy focused on stretching, strengthening, and ergonomic correction is first-line management for most musculoskeletal causes of arm pain.
 \item A carefully selected corticosteroid injection may give short-term relief for some shoulder or joint conditions, but injections near tendons can weaken them and injections for nerve entrapment have specific risks.
 \item Surgical options include carpal tunnel release, cubital tunnel decompression, rotator cuff repair, or cervical discectomy depending on the underlying pathology.
 \item Neuropathic pain agents (gabapentin, pregabalin, amitriptyline) are prescribed for radicular or neuropathic arm pain.
\end{itemize}

\section*{Self-Care and Home Management}
\begin{itemize}
 \item Protect the arm from activities that clearly worsen pain, but avoid prolonged complete rest; gentle movement within comfort can reduce stiffness.
 \item Apply ice packs for 15–20 minutes several times daily during the acute phase (first 48 hours) to reduce inflammation.
 \item Consider acetaminophen or an over-the-counter NSAID only if safe for you; NSAIDs may be unsuitable with kidney disease, ulcers, anticoagulants, heart disease, pregnancy, or certain medicines.
 \item Use ergonomic adjustments at work (keyboard tray, chair height, monitor position) to reduce strain on the arm and shoulder.
\end{itemize}

\section*{References}
\begin{thebibliography}{99}
\bibitem{arm_pain:ref1} Katz JN. ``Carpal tunnel syndrome.'' \textit{N Engl J Med}. 2020;382(8):754-762.
\bibitem{arm_pain:ref2} Whittle S, Buchbinder R. ``In the clinic: rotator cuff disease.'' \textit{Ann Intern Med}. 2015;162(1):ITC1-ITC15.
\bibitem{arm_pain:ref3} Sanders RJ, Hammond SL, et al. ``Thoracic outlet syndrome: a review.'' \textit{Neurologist}. 2008;14(6):365-373.
\bibitem{arm_pain:ref4} DeLee JC, Drez D, et al. \textit{DeLee \& Drez\'s Orthopaedic Sports Medicine}, 5th ed. Elsevier, 2022.
\bibitem{arm_pain:ref5} UpToDate. ``Evaluation of the adult with arm pain.'' Wolters Kluwer, 2023.
\bibitem{arm_pain:ref6} American College of Radiology. ACR Appropriateness Criteria: cervical pain or cervical radiculopathy. Revised 2024.
\bibitem{arm_pain:ref7} American Academy of Orthopaedic Surgeons. Management of rotator cuff injuries clinical practice guideline. Updated 2025.
\end{thebibliography}
